\begin{table}[htbp]
\centering
\small
\renewcommand{\arraystretch}{2}

\caption{Comparison of common nanopore surface modification strategies. Adapted from Eggenberger \emph{et al.}\cite{Eggenberger2019} and modified for the design requirements considered in this dissertation.}
\label{tab:SurfaceModificationComparison}

\begin{tabular}{lcccccc}
\toprule
\makecell[l]{Method} &
\makecell[c]{Phase} &
\makecell[c]{Surface\\bonding} &
\makecell[c]{Positive at\\neutral pH} &
\makecell[c]{Diameter\\control} &
\makecell[c]{Organic\\functionality} \\
\midrule

Fluid lipid coatings\cite{Yusko2011,Yusko2017,Eggenberger2019b} &
L/C &
Noncovalent &
\Yes &
\No &
\No \\

Surfactants\cite{xie2009,gangula2010} &
L/C &
Noncovalent &
\Yes &
\No &
\YesM \\

Polyelectrolyte coatings\cite{alem2007,ali2010,lepoitevin2016,ma2018,umehara2006} &
L &
Noncovalent &
\YesP &
\Yes &
\YesM \\

Self-assembled monolayers\cite{wanunu2007,nilsson2006,anderson2013,ananth2018,Sypabekova2023,ulman1996} &
L/V &
Covalent$^{*}$ &
\YesM &
\No &
\Yes \\

ALD &
V &
Covalent &
\YesM &
\YesP &
\No \\

MLD &
V &
Covalent &
\No &
\YesP &
\Yes \\

\bottomrule
\end{tabular}

\vspace{0.5em}

\begin{minipage}{0.95\textwidth}
\footnotesize
\textbf{Phase:} L = liquid phase, V = vapor phase, C = colloid.
\\
%\textbf{Covalent:} \Yes = covalently bonded, N = non-covalent.
%\\
\textbf{Scoring:} \textbf{Y$^+$}: excellent or plentiful; \textbf{\Yes}: at least some suitable options; \textbf{Y$^-$}: difficult, uncommon, 
or dubiously suitable; \textbf{N}: nominally impossible or unsuitable for hybrid nanopore formation
\\
$^{*}$Acid-base condensation (as in silanization) is considered a covalent bond, but thiol-on-gold SAMs may be considered noncovalent depending on the environment.
\end{minipage}
\renewcommand{\arraystretch}{1.0}

\end{table}